\documentclass[11pt,a4paper]{article}

%% ── Encoding & fonts ──────────────────────────────────────────────────────────
\usepackage[utf8]{inputenc}
\usepackage[T1]{fontenc}
\usepackage{lmodern}

%% ── Core packages ────────────────────────────────────────────────────────────
\usepackage{amsmath,amssymb,amsfonts}
\usepackage{mathrsfs}
\usepackage{bm}
\usepackage{geometry}
\usepackage{graphicx}
\usepackage{xcolor}
\usepackage{booktabs}
\usepackage{longtable}
\usepackage{array}
\usepackage{enumitem}
\usepackage{listings}
\usepackage{fancyhdr}
\usepackage{titlesec}
\usepackage{hyperref}
\usepackage{microtype}
\usepackage{float}
\usepackage{caption}
\usepackage{mdframed}
\usepackage{setspace}

%% ── Page geometry ────────────────────────────────────────────────────────────
\geometry{
  left=2.5cm, right=2.5cm,
  top=3.0cm,  bottom=2.8cm,
  headheight=15pt
}

%% ── Hyperref ─────────────────────────────────────────────────────────────────
\hypersetup{
  colorlinks=true,
  linkcolor=black,
  citecolor=black,
  urlcolor=black,
  pdftitle={Unified Identity, Information, and Scale-Coalescence Formulation},
  pdfauthor={Theoretical Division}
}

%% ── Listings style ───────────────────────────────────────────────────────────
\definecolor{codebg}{gray}{0.93}
\lstset{
  basicstyle=\ttfamily\small,
  backgroundcolor=\color{codebg},
  frame=single,
  framerule=0.4pt,
  rulecolor=\color{black!40},
  breaklines=true,
  breakatwhitespace=true,
  tabsize=2,
  showstringspaces=false,
  captionpos=b,
  xleftmargin=6pt,
  xrightmargin=6pt
}

%% ── Section title styling ────────────────────────────────────────────────────
\titleformat{\section}{\normalfont\large\bfseries}{\thesection.}{0.6em}{}[\vspace{-2pt}\rule{\textwidth}{0.4pt}]
\titleformat{\subsection}{\normalfont\normalsize\bfseries}{\thesubsection}{0.5em}{}
\titleformat{\subsubsection}{\normalfont\normalsize\itshape}{\thesubsubsection}{0.5em}{}

%% ── Header / footer (body pages) ─────────────────────────────────────────────
\pagestyle{fancy}
\fancyhf{}
\lhead{\footnotesize\textsc{UIISCF} \textbar{} \textit{Identity, Information, and Scale-Coalescence}}
\rhead{\footnotesize Rev.~1.0}
\lfoot{\footnotesize\today}
\rfoot{\footnotesize\thepage}
\renewcommand{\headrulewidth}{0.4pt}
\renewcommand{\footrulewidth}{0.3pt}

%% ── Notation macros ──────────────────────────────────────────────────────────
\newcommand{\Pid}{\mathscr{P}}
\newcommand{\Rec}{\Xi}
\newcommand{\Cond}{\mho}
\newcommand{\Quota}{\mathcal{Q}}
\newcommand{\Qvec}{\mathbf{Q}}
\newcommand{\avec}{\mathbf{a}}
\newcommand{\bvec}{\mathbf{b}}
\newcommand{\rvec}{\mathbf{r}}
\newcommand{\vvec}{\mathbf{v}}
\newcommand{\Fvec}{\mathbf{F}}
\newcommand{\norm}[1]{\left\lVert #1 \right\rVert}
\newcommand{\kB}{k_{\mathrm{B}}}
\newcommand{\eV}{\,\mathrm{eV}}
\newcommand{\ang}{\,\text{\AA}}
\newcommand{\fs}{\,\mathrm{fs}}
\newcommand{\amu}{\,\mathrm{amu}}
\newcommand{\lbond}{\ell_{\mathrm{bond}}}

%% ── SM's Law interjection environment ────────────────────────────────────────
\newmdenv[
  topline=true,
  bottomline=true,
  rightline=false,
  leftline=false,
  linewidth=1.2pt,
  linecolor=black,
  innertopmargin=8pt,
  innerbottommargin=8pt,
  innerleftmargin=0pt,
  innerrightmargin=0pt,
  skipabove=12pt,
  skipbelow=12pt
]{smlaw}

%% ── Captioning ───────────────────────────────────────────────────────────────
\captionsetup{font=small, labelfont=bf}

%% ═════════════════════════════════════════════════════════════════════════════
\begin{document}
%% ═════════════════════════════════════════════════════════════════════════════

%%─────────────────────────────────────────────────────────────────────────────
%% TITLE PAGE
%%─────────────────────────────────────────────────────────────────────────────
\thispagestyle{empty}

\vspace*{3cm}

\begin{center}
  \noindent\rule{\textwidth}{1.2pt}
  \vspace{1.2em}

  {\fontsize{22}{28}\selectfont\bfseries
    Unified Identity, Information,\\[0.3em]
    and Scale-Coalescence Formulation
  }

  \vspace{0.9em}
  \noindent\rule{\textwidth}{0.5pt}
  \vspace{0.9em}

  {\large\itshape
    A Theory of Particle Identity, Information Propagation,\\[0.2em]
    and Emergent Structure Across Physical Scales
  }

  \vspace{2.4em}

  {\normalsize
    \textsc{Theoretical Division}\\[0.5em]
    \textit{Particle Physics and Molecular Simulation Research Group}
  }

  \vspace{2em}

  {\small
    \begin{tabular}{rl}
      \textbf{Document identifier:} & UIISCF-1.0 \\[0.3em]
      \textbf{Revision:}            & 1.0 \\[0.3em]
      \textbf{Date:}                & \today \\[0.3em]
      \textbf{Status:}              & Theoretical specification \\
    \end{tabular}
  }

  \vspace{3cm}

  \noindent\rule{\textwidth}{0.5pt}
  \vspace{0.9em}

  \begin{minipage}{0.82\textwidth}
  \small\textbf{Abstract.}\quad
  This document formalises the Unified Identity, Information, and
  Scale-Coalescence Formulation (UIISCF): a theoretical framework that
  unifies classical molecular dynamics, subatomic particle identity,
  information accounting, and emergent structure under a single symbolic
  layer.  The formulation is grounded in measured data from high-energy
  collider experiments and calibrated against the established hadron spectrum,
  nuclear binding energies, and atomic ionisation data.  Four primitives
  organise the theory: the identity-bearing particle $\Pid_i$, the recoverable
  information $\Rec_i$, the interaction conductance $\Cond_{ij}$, and the
  active scale-identity quota $\Quota_S$.  Two governing laws constrain all
  interactions and emergence events.  Nineteen sections cover the complete
  theoretical and implementation specification, from global state representation
  through detector reconstruction, solver comparison, and version gating.
  \end{minipage}

  \vspace{0.9em}
  \noindent\rule{\textwidth}{1.2pt}
\end{center}

\newpage

%%─────────────────────────────────────────────────────────────────────────────
%% TABLE OF CONTENTS
%%─────────────────────────────────────────────────────────────────────────────
\tableofcontents
\newpage

%%═══════════════════════════════════════════════════════════════════════════════
\section{Motivation and Theoretical Basis}\label{sec:motivation}
%%═══════════════════════════════════════════════════════════════════════════════

\subsection{The Problem with Particle Identity}

In orthodox quantum field theory, a particle is a quantised excitation of a
field --- an eigenstate of the particle-number operator.  This description is
mathematically complete but operationally opaque.  At a particle collider,
no particle is ever directly observed: the detector records energy deposits,
timing signals, and ionisation trails.  What experimenters call a ``proton''
or an ``electron'' is a reconstruction: a chain of inference from raw detector
hits through track-fitting algorithms, calorimeter clustering, and
vertex-finding to an assigned identity label and four-momentum.  The particle,
in this sense, is the output of an information-processing pipeline, not a
primitive input to it.

This observation is not merely philosophical.  The reconstruction efficiency,
the mis-identification rate, the momentum resolution --- all of these are
information-theoretic quantities.  When a $\pi^+$ and a $\pi^-$ are produced
in the same collision, the question of which track belongs to which pion is a
question about recoverable information, not about the ``true'' identities of
the particles in any absolute sense.  The identities are exactly as
distinguishable as the data permits them to be.

The Unified Identity, Information, and Scale-Coalescence Formulation takes
this observation seriously and elevates it to a design principle.  Identity
is not a God-given label attached to a particle at creation; it is a quantity
of distinguishability that is budgeted, transferred, and depleted as
particles interact.  The recoverable information $\Rec_i$ is the
formal expression of this budget.

\subsection{Grounding in CERN Experimental Data}

The scale ladder at the heart of UIISCF is not arbitrary.  Each rung
corresponds to a level of matter at which experimental data is available to
constrain the parameters.

\textbf{Scale 0 --- Quarks and the Cornell potential.}
The confining potential between quarks,
$E_{\mathrm{conf}}(r) = \kappa r + \alpha_s / r$, was not invented for this
formalism.  It was extracted from quarkonium spectroscopy data: the $J/\psi$,
$\Upsilon$, and charmonium excited-state mass splittings measured at SLAC,
CERN, and the Tevatron constrain $\kappa \approx 0.18\,\mathrm{GeV}^2$ and
$\alpha_s \approx 0.12$--$0.30$ at hadronic scales.  The colour-averaged
interaction (CAI) model used here is a direct approximation to the effective
quark-antiquark potential inferred from these spectra, not a first-principles
QCD calculation.  The \texttt{full\_qcd = false} flag is mandatory precisely
because the distinction matters.

\textbf{Scale 1 --- Hadron formation and nuclear binding.}
The bonding threshold $\lbond^{(1)} \approx 10^{-4}\ang$ and the birth
criterion $m \geq 2$ (meson) or $m \geq 3$ (baryon) are set by the measured
hadron sizes from electron--proton scattering at HERA and the proton charge
radius from precision hydrogen spectroscopy.  The baryon asymmetry in
scale-1 coalescence --- that three quarks form a baryon more readily than two
--- is a direct consequence of the colour-neutrality constraint on the
triplet state.

\textbf{Scale 2 --- Atoms and classical MD.}
At the atomic scale the formalism reduces exactly to classical molecular
dynamics.  The force fields (Lennard-Jones, Morse, CHARMM, AMBER) used at
this scale are calibrated against crystal structure data, neutron scattering
spectra, and ab initio quantum chemical calculations.  The classical MD bridge
(Section~\ref{sec:cmd}) makes this reduction exact and not approximate: no
information from the subatomic layers leaks into the classical layer unless
\texttt{info\_track} is explicitly enabled.

\textbf{Scales 3--4 --- Molecules and clusters.}
Bonding thresholds at the molecular and cluster scales are set from measured
bond lengths (crystallographic databases), dissociation energies (gas-phase
thermochemistry), and cluster stability data from mass spectrometry.

\subsection{Why These Four Primitives}

The choice of $\Pid_i$, $\Rec_i$, $\Cond_{ij}$, and $\Quota_S$ as the four
primitives of the formalism is not arbitrary.  Together they form a minimal
complete basis for describing identity, information, coupling, and emergence.

The \emph{identity-bearing particle} $\Pid_i$ is the carrier: it is the
thing that has properties, occupies space, and participates in interactions.
Without an identity carrier, there is nothing to track.  The calligraphic
$\mathscr{P}$ notation is chosen to distinguish it from the Hamiltonian
operator $\hat{H}$ or the momentum $p$, with which it must coexist in the
same equations.

The \emph{recoverable information} $\Rec_i$ quantifies how distinguishable
$\Pid_i$ is from its environment.  This is the quantity that decreases when
particles interact (SM's First Law) and is partially inherited when new
identities form (SM's Second Law).  Using the Greek capital $\Xi$ --- a letter
with no strong prior claim in this domain --- deliberately avoids confusion
with entropy $S$, energy $E$, or the wavefunction $\Psi$.  The unit choice
(nats) is deliberate: nats are the natural unit of information in continuous
distributions, compatible with the thermodynamic entropy in the same units
via $S = k_{\mathrm{B}} \Rec$.

The \emph{interaction conductance} $\Cond_{ij}$ encodes the rate at which
information flows between $\Pid_i$ and $\Pid_j$ when they interact.  The
siemens-derived symbol $\mho$ is a deliberate analogy: just as electrical
conductance governs charge flow between two nodes, interaction conductance
governs information flow between two identities.  A high-conductance pair
exchanges information rapidly; a zero-conductance pair does not interact.

The \emph{scale-identity quota} $\Quota_S$ counts how many active identities
exist at each scale.  It is an integer quantity because identity is discrete:
you cannot have half a proton.  Its dynamics (SM's Second Law) describe
emergence --- the birth of higher-scale identities from lower-scale
coalescence --- in a way that is countable, verifiable, and conserved within
the constraints of baryon number and lepton number.

\subsection{A More Intuitive Representation of Matter}

The standard approach to multi-scale simulation runs separate models at
separate scales and couples them at boundaries: a quantum mechanics / molecular
mechanics (QM/MM) boundary, or a coarse-grained / fine-grained boundary.
These boundaries are artificial and generate artefacts: force discontinuities,
energy drift, and ambiguity about which particles belong to which region.

UIISCF removes the explicit boundary by treating scale as a dynamic property
of the identity, not a property of the simulation domain.  A carbon atom is a
scale-2 identity.  If sufficient energy is deposited, it can be ionised (losing
recoverable information to the environment) or, in extreme conditions,
resolved into its nuclear constituents (a scale-1 event, only reachable with
the subatomic channels active).  The same state vector $\Omega(t)$ holds the
atom regardless of which channels are active; the active channels determine
which interactions are resolved and which are approximated.

This architecture mirrors how matter actually works: the nucleus is always
there inside the atom, but at chemical energies its internal degrees of
freedom are frozen out.  UIISCF makes this freezing-out explicit through
version gating and scale-channel weights, rather than hiding it inside a
force field.  The result is a representation of matter that can be understood
from the bottom up, starting with quarks and building to clusters, rather
than bolted together from separate codes that happen to share a boundary.

%%═══════════════════════════════════════════════════════════════════════════════
\section{Notation Layer}\label{sec:notation}
%%═══════════════════════════════════════════════════════════════════════════════

Four primitive symbols are introduced.  All subsequent formalism is expressed
in terms of these primitives and the standard MD state.  The motivation for
each choice is discussed in Section~\ref{sec:motivation}; here the symbols
are defined precisely.

\begin{table}[H]
\centering
\caption{UIISCF primitive notation.}
\label{tab:notation}
\begin{tabular}{@{}clll@{}}
\toprule
Symbol & \LaTeX\ macro & Meaning & Units / Domain \\
\midrule
$\Pid_i$ & \texttt{\textbackslash mathscr\{P\}\_i} &
  Identity-bearing particle at scale $n$ &
  Dimensionless label \\
$\Rec_i$ & \texttt{\textbackslash Xi\_i} &
  Recoverable distinguishability (information) of $\Pid_i$ &
  nats \\
$\Cond_{ij}$ & \texttt{\textbackslash mho\_\{ij\}} &
  Interaction conductance between $\Pid_i$ and $\Pid_j$ &
  $[\mathrm{nats\,fs^{-1}}]$ \\
$\Quota_S$ & \texttt{\textbackslash mathcal\{Q\}\_S} &
  Active scale-identity quota at scale $S$ &
  $\mathbb{Z}_{\geq 0}$ \\
\bottomrule
\end{tabular}
\end{table}

\noindent
The symbol $n$ denotes scale level: $n=0$ (quark), $n=1$ (hadron/nucleus),
$n=2$ (atom), $n=3$ (molecule), and so on.  Scale levels are open-ended;
an implementation gates which levels are active via version constraints
(Section~\ref{sec:version}).

\bigskip

The two governing laws are stated here in full and recur as interjections
throughout the formalism wherever they constrain a derivation.

\begin{smlaw}
\textbf{First Law of Identity --- Conservation of Distinguishability.}\\[4pt]
Every pairwise interaction redistributes or consumes recoverable information.
No interaction is information-neutral.  Formally:
\[
  \Pid_i \leftrightarrow \Pid_j
  \;\implies\;
  \Delta\Rec_{ij} < 0 \;\;\text{or}\;\; \Delta S_{ij} > 0 ,
\]
where $\Delta\Rec_{ij} = \Rec_i' + \Rec_j' - \Rec_i - \Rec_j$ is the change
in total recoverable information and $\Delta S_{ij}$ is the entropy produced
by the interaction.  Equality $\Delta\Rec_{ij}=0,\;\Delta S_{ij}=0$ is
forbidden for any physical interaction.
\end{smlaw}

\begin{smlaw}
\textbf{Second Law of Identity --- Scale-Coalescence Birth Criterion.}\\[4pt]
When $m \geq 3$ lower-scale identities satisfy mutual proximity and bond
maturity, a new higher-scale identity is born:
\[
  m \geq 3,\quad
  d_{ij}^{(n)} \leq \lbond^{(n)},\quad
  B_n\!\left(C_n\right) \geq \tau_B
  \;\implies\;
  \mathrm{Exist}\!\left(\Pid_*^{(n+1)}\right) = 1,\quad
  \Delta\Quota_S = +1 ,
\]
where $d_{ij}^{(n)}$ is the pairwise separation at scale $n$,
$\lbond^{(n)}$ is the scale-$n$ bonding length threshold,
$B_n(C_n)$ is the bond maturity functional evaluated on bond count $C_n$,
and $\tau_B$ is the maturity threshold.  The new identity $\Pid_*^{(n+1)}$
inherits summed mass, net charge, and a fresh $\Rec$ budget.
\end{smlaw}

%%═══════════════════════════════════════════════════════════════════════════════
\section{Global State}\label{sec:global}
%%═══════════════════════════════════════════════════════════════════════════════

The simulation state at time $t$ is the tuple
\[
  \Omega(t) \;=\;
  \bigl(
    \mathbf{R}(t),\;
    \mathbf{V}(t),\;
    \mathbf{F}(t),\;
    \mathbf{M},\;
    \mathbf{Z},\;
    \boldsymbol{\Rec}(t),\;
    \boldsymbol{\Cond}(t),\;
    \Quota_S(t),\;
    \mathbf{H}(t),\;
    E_{\mathrm{tot}}(t),\;
    t,\;
    \Delta t
  \bigr) ,
\]
where $\mathbf{R}$, $\mathbf{V}$, $\mathbf{F} \in \mathbb{R}^{N \times 3}$
are position, velocity, and force arrays over $N$ particles;
$\mathbf{M} \in \mathbb{R}^N$ is the mass vector;
$\mathbf{Z} \in \mathbb{Z}^N$ is the atomic-number / particle-species vector;
$\boldsymbol{\Rec}(t) \in \mathbb{R}^N$ is the per-particle recoverable
information vector;
$\boldsymbol{\Cond}(t)$ is the sparse conductance matrix (non-zero only for
bonded or proximate pairs);
$\Quota_S(t) \in \mathbb{Z}_{\geq 0}$ is the active identity count at each
scale;
$\mathbf{H}(t)$ is the $3\times3$ cell matrix;
$E_{\mathrm{tot}}$ is total energy; and $\Delta t$ is the current timestep.

\noindent
Compared to classical MD, the state $\Omega$ carries two new blocks:
$\boldsymbol{\Rec}$ and $\boldsymbol{\Cond}$.  These blocks are dormant
(zero-cost) when information tracking is disabled and active only when
explicitly requested.  This design ensures full backward compatibility with
existing MD workflows: a simulation that never writes to these blocks
behaves identically to a classical MD run.

\subsection{State-Vector Projection}

The universal state vector $\Qvec$ projects onto $\Omega$ via:
\[
  \Qvec(t) \;=\;
  \bigl[\,\mathbf{R}^\top \;\big|\; \mathbf{V}^\top \;\big|\;
         \boldsymbol{\Rec}^\top \;\big|\; \mathrm{vec}(\boldsymbol{\Cond}) \,\bigr]^\top .
\]
Every observable (pressure, temperature, spectra, identity counts) is a
projection $\mathcal{P}_\alpha(\Qvec)$ for some projection operator
$\mathcal{P}_\alpha$.  The identity-information block
$[\boldsymbol{\Rec}^\top \mid \mathrm{vec}(\boldsymbol{\Cond})]$ is new in
UIISCF and was absent from the classical MD projection.

\subsection{Initialisation Invariants}

At $t=0$ the following must hold:
\begin{enumerate}[leftmargin=1.5em]
  \item $\Rec_i(0) = \Rec_i^{\mathrm{birth}}$ for every particle, set from
    the species registry.
  \item $\Cond_{ij}(0) = 0$ for all pairs unless an explicit bond is declared
    in the input block.
  \item $\Quota_S(0) = N$ (all particles are scale-$n$ identities).
  \item $\mathbf{H}(0)$ is orthonormal and volume-positive:
    $\det(\mathbf{H}) > 0$.
\end{enumerate}


%%═══════════════════════════════════════════════════════════════════════════════
\section{Scale Identity Ladder}\label{sec:ladder}
%%═══════════════════════════════════════════════════════════════════════════════

The scale ladder is the structural backbone of UIISCF.  Rather than
committing to a single level of description, the formalism assigns each
identity to a rung $n$ and allows transitions between rungs through
coalescence (upward) and dissociation (downward).  The parameters at each
rung are not freely adjustable: they are anchored to measured data.  The
bonding lengths come from scattering cross-sections and structural databases;
the birth thresholds come from nuclear and molecular stability measurements;
the maturity exponents are calibrated so that the formation rate of common
hadrons and small molecules reproduces known reaction rates at the relevant
temperatures.

The reason for requiring $m \geq 3$ as the minimum coalescence count
(rather than two, as in a naive bonding model) is both physical and
computational.  Physically, two-body encounters rarely produce stable bound
states without a third body to carry away excess momentum: this is the
standard three-body recombination argument in atomic and nuclear physics.
Computationally, requiring a minimum of three participants prevents spurious
coalescence events from cascading through the identity ladder during periods
of high-density fluctuation.

Particles are organised on a discrete scale ladder.  Each rung $n$ has its own
bonding threshold, maturity functional, and information budget.

\begin{table}[H]
\centering
\caption{Scale identity ladder rungs.  Thresholds are calibrated from
  experimental data as described in Section~\ref{sec:motivation}.}
\label{tab:ladder}
\begin{tabular}{@{}clccc@{}}
\toprule
Scale $n$ & Entity type & $\lbond^{(n)}$ (\AA) &
  $\tau_B$ & Birth criterion (min.\ $m$) \\
\midrule
0 & Quark & $10^{-5}$ & 0.90 & 3 (colour-neutral triplet) \\
1 & Hadron / nucleus & $10^{-4}$ & 0.85 & 2 (meson) or 3 (baryon) \\
2 & Atom & 2.5 & 0.70 & 2 \\
3 & Molecule & 5.0 & 0.65 & 3 \\
4 & Cluster & 15.0 & 0.60 & 5 \\
\bottomrule
\end{tabular}
\end{table}

\subsection{Bond Maturity Functional}

The bond maturity functional is chosen to be a saturating exponential:
\[
  B_n(C_n) \;=\; 1 - \exp\!\left(-\lambda_n C_n\right),
\]
where $C_n$ is the cumulative bond-event count at scale $n$ and
$\lambda_n > 0$ is the maturity rate parameter.  This functional form is
chosen deliberately.  A linear maturity criterion would allow coalescence
to trigger on the first interaction, which is too aggressive.  A hard
threshold on $C_n$ would produce discontinuous birth rates.  The saturating
exponential is the simplest functional that is continuous, monotonically
increasing, bounded in $[0,1)$, and parameterised by a single rate constant
whose physical interpretation is straightforward: $1/\lambda_n$ is the
characteristic interaction count for maturation.  Default values:
$\lambda_0 = 5,\; \lambda_1 = 3,\; \lambda_2 = 1.5,\; \lambda_3 = 0.8$.

\subsection{Information Budget at Birth}

When the Second Law of Identity fires and $\Pid_*^{(n+1)}$ is born:
\[
  \Rec_*^{(n+1)} \;=\;
  \eta \sum_{i \in \mathcal{C}} \Rec_i^{(n)}
  \;-\; \Delta S_{\mathrm{coal}} ,
\]
where $\mathcal{C}$ is the coalescing set, $\eta \in (0,1]$ is the
information-transfer efficiency (default $\eta = 0.95$), and
$\Delta S_{\mathrm{coal}} > 0$ is the coalescence entropy cost.  The deficit
$\sum_i \Rec_i - \Rec_*/\eta$ is written to the entropy ledger.

\begin{smlaw}
\textbf{First Law of Identity} applies at coalescence: the information loss
$\sum_i \Rec_i^{(n)} - \Rec_*^{(n+1)} > 0$ must be strictly positive.  A
coalescence that conserves or increases total recoverable information is
rejected with error \texttt{ERR\_COAL\_INFO\_NONZERO}.
\end{smlaw}

%%═══════════════════════════════════════════════════════════════════════════════
\section{Classical MD Bridge}\label{sec:cmd}
%%═══════════════════════════════════════════════════════════════════════════════

For scale-2 (atomic) identities the UIISCF reduces exactly to standard
molecular dynamics.  This is not a coincidence or a special case; it is an
explicit design requirement.  The classical MD literature represents decades
of validated force fields, thermostat algorithms, barostat algorithms, and
long-range solvers.  Any formalism that introduces a subatomic layer while
breaking the classical layer loses the trust of the community that already
uses classical MD.  The bridge is therefore designed so that activating
information tracking adds exactly two new ODEs per particle ($\dot{\Rec}_i$
and $\dot{\Cond}_{ij}$) without modifying the existing force, position, or
velocity equations.

\subsection{Equations of Motion}

\[
  m_i \ddot{\rvec}_i \;=\; \Fvec_i \;=\; -\nabla_{\rvec_i} U ,
\]
\[
  U \;=\; \sum_{i<j} \phi(r_{ij}) \;+\;
          \sum_{\mathrm{bonds}} V_{\mathrm{bond}} \;+\;
          \sum_{\mathrm{angles}} V_{\mathrm{angle}} \;+\;
          \sum_{\mathrm{dihedrals}} V_{\mathrm{dihed}} .
\]
Periodic boundary conditions apply via the minimum-image convention using
cell matrix $\mathbf{H}$:
\[
  \rvec_{ij}^{\mathrm{MIC}} \;=\;
  \mathbf{H}\!\left(
    \mathbf{s}_{ij} - \mathrm{round}(\mathbf{s}_{ij})
  \right),
  \quad
  \mathbf{s}_{ij} = \mathbf{H}^{-1}(\rvec_i - \rvec_j) .
\]

\subsection{Information Coupling to MD}

At scale 2 the recoverable information $\Rec_i$ is a passive bookkeeping
quantity computed alongside positions.  For every force evaluation step:
\[
  \dot{\Rec}_i \;=\;
  -\sum_{j \neq i} \Cond_{ij} \,\bigl(\Rec_i - \Rec_j\bigr)
  \;-\; \gamma_i \Rec_i ,
\]
where $\gamma_i \geq 0$ is a species-specific information dissipation rate and
the conductance $\Cond_{ij}$ is proportional to the pair interaction strength.
This equation has the structure of a heat diffusion equation on the interaction
graph: information flows from high-$\Rec$ particles to low-$\Rec$ particles
along edges weighted by conductance, and dissipates at rate $\gamma_i$.
The analogy with heat conduction is intentional and makes the physical
interpretation transparent.

\begin{lstlisting}[language={}, caption={Classical MD configuration block (reference scripting format).}]
[run]
integrator   = "velocity_verlet"
dt           = 0.5          # fs
steps        = 200000
thermostat   = "nose_hoover"
barostat     = "parrinello_rahman"
info_track   = true         # enable Xi / mho bookkeeping
info_dissip  = 0.001        # gamma_i default (fs^-1)
\end{lstlisting}

%%═══════════════════════════════════════════════════════════════════════════════
\section{Matrix-Force Policy}\label{sec:matforce}
%%═══════════════════════════════════════════════════════════════════════════════

Classical MD codes hard-code force-computation loops: one loop for pair
potentials, one for bonds, one for angles, and so on.  Adding a subatomic
force channel in this paradigm requires modifying the innermost loops, which
risks breaking validated classical behaviour and makes it difficult to gate
subatomic contributions by scale or version.

The matrix-force policy solves this by routing all force computation through
a weighted dispatch table.  Each force channel contributes a force vector
$\Fvec_i^{(\alpha)}$; the total force is a weighted sum.  Channel weights
are set to zero by default for all subatomic channels and raised to 1.0 only
when the corresponding scale is active.  This architecture means that adding
a new channel --- colour-averaged interaction, gluon proxy, weak proxy ---
never modifies an existing channel, and can always be validated in isolation
before being combined.  The classical path is preserved exactly: with all
subatomic weights at zero, the dispatch table produces the identical result
to the original hard-coded loops.

\subsection{Force Dispatch Table}

\[
  \Fvec_i \;=\; \sum_{\alpha} w_\alpha \, \Fvec_i^{(\alpha)} ,
\]
where $\alpha$ indexes active force channels and $w_\alpha$ is a channel
weight gated by the version and active-scale flags.

\begin{table}[H]
\centering
\caption{Force channel dispatch table.}
\label{tab:dispatch}
\begin{tabular}{@{}lllc@{}}
\toprule
Channel $\alpha$ & Description & Default $w_\alpha$ & Min.\ version \\
\midrule
\texttt{pair}    & Pair potential (LJ, Morse, \ldots) & 1.0 & v5.1.4 \\
\texttt{bond}    & Harmonic/FENE bond & 1.0 & v5.1.4 \\
\texttt{angle}   & Valence angle & 1.0 & v5.1.4 \\
\texttt{dihedral}& Dihedral torsion & 1.0 & v5.1.4 \\
\texttt{coulomb} & Long-range Coulomb (PPPM) & 1.0 & v5.1.4 \\
\texttt{cai}     & Colour-averaged interaction (subatomic) & 0.0 & v5.1.5 \\
\texttt{gluon}   & Cornell proxy (effective QCD) & 0.0 & v5.1.5 \\
\bottomrule
\end{tabular}
\end{table}

\subsection{Policy Enforcement}

The engine validates the weight vector before every run:
\begin{enumerate}[leftmargin=1.5em]
  \item $\sum_\alpha w_\alpha > 0$ (at least one active channel).
  \item Subatomic channels (\texttt{cai}, \texttt{gluon}) require
    \texttt{subatomic\_mode = true} and kernel $\geq$ v5.1.5.
  \item If \texttt{gluon} weight $> 0$ then \texttt{full\_qcd = false}
    \emph{must} be explicitly set.  Omission raises
    \texttt{ERR\_QCD\_FLAG\_MISSING}.
\end{enumerate}

\begin{smlaw}
\textbf{First Law of Identity} constrains channel coupling: enabling a new
force channel always changes $\Rec$ accounting.  Adding \texttt{cai} or
\texttt{gluon} without enabling \texttt{info\_track} is forbidden; the
implementation will reject the configuration with
\texttt{ERR\_INFO\_TRACK\_REQUIRED}.
\end{smlaw}

%%═══════════════════════════════════════════════════════════════════════════════
\section{Per-Atom Event-Energy State}\label{sec:event}
%%═══════════════════════════════════════════════════════════════════════════════

Continuous force integration handles the smooth evolution of positions and
velocities but is not the right tool for discontinuous events: bond formation,
bond breakage, coalescence, annihilation.  These events change the topology
of the interaction graph instantaneously and deposit or release discrete
amounts of energy.  Tracking them requires a dedicated event-energy register
per particle, separate from the running kinetic and potential energy.

Separating event energy from kinetic and potential energy also provides a
natural diagnostic: the event-energy trace is a record of every discrete
transition that occurred, in time order, with position, energy, and
information-change metadata.  This trace can be replayed, audited, and
compared across solver runs independently of the continuous trajectory.

Each particle carries:
\[
  \mathcal{E}_i(t) \;=\;
  \bigl(
    E_i^{\mathrm{kin}},\;
    E_i^{\mathrm{pot}},\;
    E_i^{\mathrm{event}},\;
    \Rec_i,\;
    n_i^{\mathrm{evt}}
  \bigr) ,
\]
where $E_i^{\mathrm{event}}$ is the cumulative energy from discrete events and
$n_i^{\mathrm{evt}}$ is the event counter.

\subsection{Event Types}

\begin{table}[H]
\centering
\caption{Discrete event types and their energy/information effects.}
\label{tab:events}
\begin{tabular}{@{}lllll@{}}
\toprule
Event type & Code & $\Delta E^{\mathrm{event}}$ & $\Delta\Rec$ & $\Delta\Quota_S$ \\
\midrule
Bond formation & \texttt{BOND\_FORM} & $-E_{\mathrm{bond}}$ & $< 0$ & 0 \\
Bond break     & \texttt{BOND\_BREAK} & $+E_{\mathrm{bond}}$ & $< 0$ & 0 \\
Coalescence    & \texttt{COAL} & $-\Delta E_{\mathrm{coal}}$ & $\leq -\Delta S_{\mathrm{coal}}$ & $+1$ \\
Dissociation   & \texttt{DISSOC} & $+\Delta E_{\mathrm{dissoc}}$ & $< 0$ & $-1$ \\
Annihilation   & \texttt{ANNIH} & $+2m_ic^2$ & $-\Rec_i - \Rec_{\bar\imath}$ & $-2$ \\
Excitation     & \texttt{EXCITE} & $+E_\gamma$ & $< 0$ & 0 \\
De-excitation  & \texttt{DEEXCITE} & $-E_\gamma$ & $< 0$ & 0 \\
\bottomrule
\end{tabular}
\end{table}

\subsection{Event Log Format}

Each event is appended to the event log (\texttt{.evlog} artifact) as:
\begin{lstlisting}
STEP  TYPE        ID_A   ID_B   SCALE  dE_eV        dXi
12500 BOND_FORM   1042   2107   2      -4.32100e+00  -0.0023
12501 COAL        0133   0134   1      -1.20000e+02  -1.4400
\end{lstlisting}

%%═══════════════════════════════════════════════════════════════════════════════
\section{Event Limiter}\label{sec:limiter}
%%═══════════════════════════════════════════════════════════════════════════════

Unconstrained event rates can destabilise numerical integration.  A single
timestep that generates hundreds of coalescence events deposits a large
impulsive energy into the system, causing temperature spikes and, in extreme
cases, the velocity-Verlet integrator to step a particle through a potential
barrier.  The event limiter is a numerical stability mechanism, not a physical
one: it does not claim that real matter has a maximum event rate, only that
the discrete-time integrator does.

The rate ceiling is set proportional to $N \cdot \Delta t$ so that the
ceiling scales appropriately with system size and timestep, avoiding the
pathological situation where a fine-grained timestep unexpectedly allows
more events per step than a coarse one.

\subsection{Rate Ceiling}

\[
  n_{\mathrm{evt}}^{\mathrm{step}} \;\leq\; \left\lfloor
    \lambda_{\mathrm{evt}} \cdot N \cdot \Delta t
  \right\rfloor ,
\]
where $\lambda_{\mathrm{evt}}$ is the maximum event rate density
(events per particle per fs), default $\lambda_{\mathrm{evt}} = 0.01$.
Events exceeding the ceiling in a single step are deferred to the next step
via a FIFO queue.

\subsection{Energy-Change Guard}

For any single event:
\[
  \left|\Delta E_i^{\mathrm{event}}\right| \;\leq\;
  \alpha_{\mathrm{guard}} \cdot E_i^{\mathrm{kin}} ,
\]
where $\alpha_{\mathrm{guard}} = 5.0$ by default.  Events violating this
guard are logged as \texttt{GUARD\_CLIP} and the energy change is clamped.

\begin{smlaw}
\textbf{First Law of Identity} applies to clamped events: even a clipped
energy change must satisfy $\Delta\Rec < 0$.  A clamp that would produce
$\Delta\Rec \geq 0$ is rejected entirely and logged as
\texttt{ERR\_GUARD\_INFO\_VIOLATION}.
\end{smlaw}

\begin{lstlisting}[language={}, caption={Event limiter configuration block.}]
[kernel]
event_rate_density  = 0.01    # events per particle per fs
event_energy_guard  = 5.0     # alpha_guard
event_queue_depth   = 1000    # max deferred events
event_log_path      = "sim.evlog"
\end{lstlisting}

%%═══════════════════════════════════════════════════════════════════════════════
\section{Formation-Defect Bridge}\label{sec:formation}
%%═══════════════════════════════════════════════════════════════════════════════

Coalescence and dissociation leave traces in the trajectory record beyond
the continuous position and velocity data.  The formation-defect bridge is
the mechanism by which these traces are captured in a dedicated formation
trajectory format and linked back to the identity ladder.

The information-loss fraction $\delta_{\mathrm{defect}}$ at dissociation
provides a physically meaningful severity classification that has no analogue
in purely energetic defect classification schemes.  A clean split
($\delta < 0.02$) indicates that the two fragments carry nearly the same
distinguishability they had before binding; the information was stored in
the interaction, not consumed.  A fragmentation event ($\delta > 0.1$)
indicates that most of the identity information of the parent was lost to
the environment --- the fragments are fundamentally different objects,
not simply pieces of the original.

The formation trajectory connects formation events to coalescence transitions
on the identity ladder.  Defect events connect to dissociation and bond-break
transitions.

\subsection{Formation Event Record}

A formation event is triggered when the Second Law of Identity fires:
\begin{lstlisting}
# formation event at step 50000
FORM  parent_ids=[1042,1043,1044]  child_id=5000  scale=2->3
      mass_sum=28.054 charge_sum=0  Xi_child=2.8500
      dS_coal=0.0247  bond_maturity=0.9200
\end{lstlisting}

\subsection{Defect Classification}

Defects are classified by the information change at dissociation:
\[
  \delta_{\mathrm{defect}} \;=\;
  \frac{\Rec_{\mathrm{parent}} -
        \sum_{k} \Rec_k^{\mathrm{fragment}}}
       {\Rec_{\mathrm{parent}}} .
\]
\begin{table}[H]
\centering
\caption{Defect classification by information loss fraction.}
\label{tab:defect}
\begin{tabular}{@{}llll@{}}
\toprule
Class & $\delta_{\mathrm{defect}}$ range & Label & Severity \\
\midrule
Clean split   & $[0, 0.02)$  & \texttt{DEFECT\_CLEAN}  & Low \\
Lossy split   & $[0.02, 0.1)$& \texttt{DEFECT\_LOSSY}  & Medium \\
Fragmentation & $[0.1, 1.0]$ & \texttt{DEFECT\_FRAG}   & High \\
\bottomrule
\end{tabular}
\end{table}

\begin{smlaw}
\textbf{Second Law of Identity} sets the lower bound on coalescence count:
$m \geq 3$.  Any two-body coalescence at scale $n \geq 2$ is promoted to a
\texttt{COAL\_PROBE} event and only finalised if a third proximate particle
joins within $\tau_{\mathrm{probe}} = 5\fs$ (configurable).
\end{smlaw}


%%═══════════════════════════════════════════════════════════════════════════════
\section{Quark Component Matrix}\label{sec:quark}
%%═══════════════════════════════════════════════════════════════════════════════

The quark sector is the point at which this formalism is most directly
accountable to measured data.  The parameters that appear here are not
fitted to a simulation objective; they are inherited from experimental
hadron physics.  The Cornell potential string tension $\kappa \approx
0.18\,\mathrm{GeV}^2$ is the value inferred from the linear Regge trajectories
of light-quark mesons as measured at CERN's NA series and the Fermilab fixed-
target experiments.  The running coupling $\alpha_s \approx 0.18$ at the
relevant hadronic scale is consistent with the world average from the PDG
compilation of $\alpha_s(m_Z) = 0.1179 \pm 0.0010$ evolved down to
$Q \sim 1\,\mathrm{GeV}$ using two-loop running.  The hyperfine coupling
$V_{\mathrm{hyp}}$ is calibrated from the $J/\psi$--$\eta_c$ and
$\Upsilon$--$\eta_b$ mass splittings.

Working backwards from CERN collision data to a tractable simulation model
reveals something useful: quark identity is less a property of the quark
itself than of the hadron it inhabits.  A quark confined in a proton is
a different object --- in terms of its effective mass, its coupling, and its
contribution to conserved quantum numbers --- from a quark in a pion or in
the deconfined quark-gluon plasma produced in heavy-ion collisions at the LHC.
The quark component matrix encodes this context-dependence explicitly through
the recoverable information column $\Rec_i^{(0)}$, which varies between
hadronic environments.

At scale $n=0$ each quark is an identity $\Pid_i^{(0)}$ carrying colour
charge $c_i \in \{R, G, B, \bar R, \bar G, \bar B\}$, flavour $f_i$, and
spin $\sigma_i = \pm\tfrac{1}{2}$.  The quark component matrix $\mathbf{Q}^{(0)}$
has one row per quark:
\[
  Q^{(0)}_{ik} \;=\;
  \begin{cases}
    m_i / m_u     & k=1 \quad\text{(mass in units of } m_u\text{)} \\
    q_i           & k=2 \quad\text{(electric charge)} \\
    c_i^{\mathrm{int}} & k=3 \quad\text{(colour integer: } 1\ldots 6\text{)} \\
    f_i^{\mathrm{int}} & k=4 \quad\text{(flavour: } u=1,d=2,s=3,c=4,b=5,t=6\text{)} \\
    \sigma_i       & k=5 \quad\text{(spin: } \pm 1\text{)} \\
    \Rec_i^{(0)}   & k=6 \quad\text{(recoverable information)} \\
  \end{cases}
\]

\subsection{Colour-Neutral Hadron Formation}

Quarks form hadrons when a colour-neutral triplet (or pair for mesons)
satisfies the Second Law of Identity at scale 0.  The colour-neutrality
condition:
\[
  \bigoplus_{i \in \mathcal{C}} c_i \;=\; \text{white} ,
\]
where $\oplus$ denotes colour addition.  Colour neutrality is checked
before applying the birth criterion.

\subsection{Colour-Averaged Interaction (CAI)}

Because full QCD is not implemented, a colour-averaged interaction (CAI)
channel approximates the strong force between quarks.  The CAI potential:
\[
  V_{\mathrm{CAI}}(r_{ij}) \;=\;
  -\frac{\alpha_s}{r_{ij}} \,\delta_{c_i \bar c_j}
  \;+\; \kappa_{\mathrm{str}} r_{ij}
  \;+\; V_{\mathrm{hyp}} ,
\]
where $\alpha_s \approx 0.12$--$0.30$, $\kappa_{\mathrm{str}} \approx
0.18\,\mathrm{GeV}^2$ (Cornell string tension), and $V_{\mathrm{hyp}}$ is
the hyperfine correction from spin-colour coupling.

\begin{lstlisting}[language={}, caption={Quark configuration block (reference implementation).}]
[quark]
enabled       = true
flavours      = ["u","d","s"]  # active flavour set
alpha_s       = 0.18
kappa_string  = 0.18           # GeV^2
hyp_coupling  = 0.02
full_qcd      = false          # MANDATORY: always set false
cai_channel   = true
info_track    = true
\end{lstlisting}

%%═══════════════════════════════════════════════════════════════════════════════
\section{Lepton Identity Matrix}\label{sec:lepton}
%%═══════════════════════════════════════════════════════════════════════════════

Leptons are fundamental-scale identities that do not participate in the strong
interaction.  They are among the most precisely measured particles in the
Standard Model: the electron mass is known to 13 significant figures, the
muon anomalous magnetic moment to 10.  This precision makes the lepton sector
an unusually clean test of identity accounting: the recoverable information of
an electron is its electric charge, mass, and spin --- three quantities that
never change in any interaction short of annihilation.  $\Rec_e$ is therefore
essentially constant until the electron annihilates, at which point it drops
to zero discontinuously.  This is the clearest example of the First Law of
Identity at work: no interaction changes the electron's $\Rec$ until the one
interaction that destroys it entirely.

Each lepton $\Pid_i^{(\ell)}$ occupies one row of the lepton identity matrix
$\mathbf{L}$:
\[
  L_{ik} \;=\;
  \begin{cases}
    m_i / m_e     & k=1 \quad\text{(mass in electron masses)} \\
    q_i           & k=2 \quad\text{(charge: } -1, 0\text{)} \\
    g_i           & k=3 \quad\text{(generation: 1, 2, 3)} \\
    \sigma_i       & k=4 \quad\text{(spin)} \\
    L_i^{\mathrm{num}} & k=5 \quad\text{(lepton number: } \pm 1\text{)} \\
    \Rec_i^{(\ell)} & k=6 \quad\text{(recoverable information)} \\
  \end{cases}
\]

\subsection{Lepton Interactions}

Leptons interact via three channels in the reference implementation:
\begin{enumerate}[leftmargin=1.5em]
  \item \textbf{Electromagnetic:} Coulomb interaction with all charged
    particles via the PPPM long-range solver.
  \item \textbf{Weak proxy:} A Yukawa-type screened potential approximating
    $W/Z$-boson exchange (range $\sim 10^{-3}\ang$, gated by
    \texttt{weak\_proxy = true}).
  \item \textbf{Electron shell binding:} At scale 2, electrons are not
    simulated as explicit particles by default; their effect enters via the
    force field.  Explicit electron tracking requires
    \texttt{explicit\_electrons = true} (v5.1.5+).
\end{enumerate}

\begin{smlaw}
\textbf{First Law of Identity} is especially active for lepton-antilepton
encounters.  An electron meeting a positron undergoes annihilation
(\texttt{ANNIH}), releasing $\Delta\Rec = -\Rec_e - \Rec_{\bar e}$ and
energy $2m_e c^2 + E_{\mathrm{kin}}$.  The information is written to the
entropy ledger irreversibly.
\end{smlaw}

%%═══════════════════════════════════════════════════════════════════════════════
\section{Boson Identity Matrix}\label{sec:boson}
%%═══════════════════════════════════════════════════════════════════════════════

Gauge bosons occupy an unusual position in the UIISCF framework.  In the
Standard Model they are the carriers of forces; in the simulation they are
either handled analytically (photons, via potential energy) or approximated
by proxy potentials (gluons, via the Cornell model; $W/Z$, via Yukawa
screening).  The decision to represent bosons as transient matrix entries
rather than persistent identities reflects a deliberate modelling choice:
at the timescales and length scales of the simulation, an exchanged virtual
photon has a lifetime many orders of magnitude shorter than the timestep.
Treating it as a persistent particle would be physically inaccurate and
computationally wasteful.

The boson identity matrix is therefore a log of mediator events, not a
roster of persistent particles.

Gauge bosons are mediator particles represented as transient entries in
the boson identity matrix $\mathbf{B}$:
\[
  B_{ik} \;=\;
  \begin{cases}
    m_i^{\mathrm{eff}} & k=1 \quad\text{(effective mass in GeV$/c^2$)} \\
    q_i               & k=2 \quad\text{(charge)} \\
    s_i               & k=3 \quad\text{(spin: 0, 1, 2)} \\
    \tau_i^{\mathrm{life}} & k=4 \quad\text{(lifetime in fs, or } \infty\text{)} \\
    \mathrm{type}_i   & k=5 \quad\text{(photon, $W^\pm$, $Z^0$, gluon, Higgs proxy)} \\
    \Rec_i^{(b)}       & k=6 \quad\text{(near zero for gauge bosons)} \\
  \end{cases}
\]

\subsection{Gluon Proxy Model}

\textbf{The gluon proxy is an effective model.  Full QCD colour dynamics are
not implemented.}  The Cornell potential governs quark confinement in place of
explicit gluon exchange:
\[
  E_{\mathrm{conf}}(r) \;=\; \kappa r \;+\; \frac{\alpha_s}{r},
  \quad
  \kappa \approx 0.18\,\mathrm{GeV}^2,
  \quad
  \alpha_s \approx 0.1\text{--}0.3 .
\]
This is a Cornell string-plus-Coulomb potential calibrated to reproduce
hadron mass splittings at low energy.  Its parameters are taken directly
from quarkonium spectroscopy, as described in Section~\ref{sec:quark}.
It is \emph{not} a perturbative QCD calculation.

\begin{smlaw}
\textbf{Mandatory disclaimer:} Any configuration using quark or boson channels
must explicitly set \texttt{full\_qcd = false}.  Omission raises
\texttt{ERR\_QCD\_FLAG\_MISSING} at parse time.  This requirement exists to
prevent Cornell-proxy results from being cited as genuine QCD predictions.
\end{smlaw}

\subsection{Photon Emission Proxy}

Photon emission (de-excitation) is handled as an event in the event log,
not as a propagating particle.  The emitted energy is $E_\gamma = \Delta E$,
recorded as \texttt{DEEXCITE}.  Photon absorption is the time-reverse.

%%═══════════════════════════════════════════════════════════════════════════════
\section{Relation-Particle Matrix}\label{sec:relation}
%%═══════════════════════════════════════════════════════════════════════════════

Classical MD tracks particles and their properties but rarely tracks the
history of individual interactions in structured form.  The relation-particle
matrix fills this gap by recording, for every currently active particle pair,
the current conductance, separation, information history, and bond type.  This
matrix is the information substrate on which the First Law operates: the
$\Delta\Rec_{ij}^{\mathrm{last}}$ column is a direct record of information
transferred at each event, and $\Delta S_{ij}^{\mathrm{cum}}$ is the running
entropy produced by the pair.

The relation-particle matrix $\mathbf{R}_{\mathrm{rel}}$ encodes pairwise
relationships.  Each row is a directed pair $(i, j)$:
\[
  (R_{\mathrm{rel}})_{(ij),k} \;=\;
  \begin{cases}
    \Cond_{ij}       & k=1 \quad\text{(interaction conductance)} \\
    d_{ij}            & k=2 \quad\text{(current separation)} \\
    \Delta\Rec_{ij}^{\mathrm{last}} & k=3 \quad\text{(information delta at last event)} \\
    \Delta S_{ij}^{\mathrm{cum}}    & k=4 \quad\text{(cumulative entropy from pair)} \\
    t_{ij}^{\mathrm{last}}          & k=5 \quad\text{(time of last event)} \\
    \mathrm{bond}_{ij}              & k=6 \quad\text{(bond type: 0=none, 1=covalent, 2=ionic,\ldots)} \\
  \end{cases}
\]

\subsection{Conductance Dynamics}

The conductance $\Cond_{ij}$ evolves between force steps:
\[
  \dot{\Cond}_{ij} \;=\;
  \alpha_{\Cond} \bigl(
    \Cond_{ij}^{\mathrm{eq}}(r_{ij}) - \Cond_{ij}
  \bigr)
  \;-\; \beta_{\Cond} \Cond_{ij} ,
\]
where $\Cond_{ij}^{\mathrm{eq}}(r)$ is the equilibrium conductance
(decreasing with $r$, zero beyond $r_{\mathrm{cut}}$), $\alpha_{\Cond}$
is the relaxation rate, and $\beta_{\Cond}$ is the natural decay rate.

%%═══════════════════════════════════════════════════════════════════════════════
\section{Proper-Time Identity Matrix}\label{sec:propertime}
%%═══════════════════════════════════════════════════════════════════════════════

At sub-relativistic velocities ($v \ll c$) proper time coincides with
simulation time and the proper-time identity matrix is diagonal:
$T_{ii} = t$.  The full relativistic extension is anticipated for a
future release and is sketched here for completeness.

\subsection{Lorentz Factor Correction}

For a particle with velocity $v_i$:
\[
  \tau_i(t) \;=\; \int_0^t
    \sqrt{1 - \frac{v_i^2(t')}{c^2}} \, dt'
  \;\approx\;
  t \left(1 - \frac{\langle v_i^2 \rangle}{2c^2}\right) .
\]
The proper-time lag $T_{ij} = \tau_i(t) - \tau_j(t)$ for $i \neq j$
becomes significant in beam-simulation scenarios where particles are
accelerated to relativistic energies: two quarks produced in the same
collision but boosted in opposite directions accumulate different proper
times and therefore age their identity information at different rates.

\subsection{Information Age}

The proper-time formalism provides a natural measure of information age:
\[
  \Rec_i^{\mathrm{age}}(\tau) \;=\;
  \Rec_i(0) \exp(-\kappa_{\mathrm{age}} \tau_i) ,
\]
where $\kappa_{\mathrm{age}}$ is an information half-life parameter, relevant
for long-lived radioactive species where the identity information of the
parent nucleus degrades as it approaches decay.

\begin{smlaw}
\textbf{First Law of Identity} in the relativistic frame: the interaction
must satisfy $\Delta\Rec_{ij} < 0$ or $\Delta S_{ij} > 0$ in the lab frame.
Lorentz boosts do not exempt an interaction from the law; information
accounting is performed in the lab frame before any boost is applied.
\end{smlaw}

%%═══════════════════════════════════════════════════════════════════════════════
\section{Dark-Sector Projection}\label{sec:dark}
%%═══════════════════════════════════════════════════════════════════════════════

The dark-sector projection is included in the UIISCF not because the nature
of dark matter is understood, but precisely because it is not.  The
information-theoretic framework provides a natural language for dark matter:
a dark-sector particle $\Pid_i^{(\mathrm{dk})}$ is an identity that carries
mass and therefore participates in gravitational clustering, but carries
zero electromagnetic charge and therefore contributes nothing to detector
hits.  Its recoverable information is sequestered: it cannot be transferred
to standard-model particles through the channels available at collider
energies.

This sequestration is the formal expression of why dark matter is dark.
It is not absent from the simulation; it is present but its $\Rec$ is not
accessible to any standard-model detector projection operator.

Dark-sector particles carry mass but do not participate in electromagnetic
or strong interactions:
\[
  D_{ik} \;=\;
  \begin{cases}
    m_i^{\mathrm{dk}}    & k=1 \\
    0                    & k=2 \quad\text{(no EM charge)} \\
    \sigma_i^{\mathrm{dk}} & k=3 \\
    \xi_i^{\mathrm{dk}}  & k=4 \quad\text{(dark coupling)} \\
    \Rec_i^{(\mathrm{dk})} & k=5 \\
  \end{cases}
\]

\subsection{Dark-Sector Force Channel}

The dark-sector force is modelled as a Yukawa potential:
\[
  V_{\mathrm{dk}}(r_{ij}) \;=\;
  -\xi_i \xi_j \, \frac{e^{-m_{\mathrm{med}} r_{ij}}}{r_{ij}} ,
\]
where $m_{\mathrm{med}}$ is the dark mediator mass.  In the limit
$m_{\mathrm{med}} \to 0$ this reduces to a gravitational-like $1/r$ potential.
The projection is explicitly \emph{phenomenological}: no claim is made
regarding correspondence with any particular dark matter model.

\begin{lstlisting}[language={}, caption={Dark-sector configuration block.}]
[dark_sector]
enabled        = false         # off by default
mediator_mass  = 1.0e-4       # GeV/c^2
coupling       = 0.001
info_track     = true
full_qcd       = false
\end{lstlisting}

%%═══════════════════════════════════════════════════════════════════════════════
\section{Event/Annihilation Matrix}\label{sec:annihilation}
%%═══════════════════════════════════════════════════════════════════════════════

Annihilation is the most dramatic event in the UIISCF framework: two
identities are extinguished and their rest mass converted to radiation.
From an information-theoretic perspective, annihilation is the maximum
possible violation of identity continuity.  The incoming particles have
well-defined identity ($\Rec_e, \Rec_{\bar e} > 0$); the outgoing photons
carry no recoverable identity information ($\Rec_\gamma \approx 0$).
The entire information budget of two particles is written to the entropy
ledger in a single step.

The annihilation matrix records this event comprehensively.  The record is
designed to be auditable after the fact: given the annihilation matrix, one
can verify that lepton number was conserved (one electron and one positron
entered), that energy-momentum was conserved (the photon energies sum to
$2m_e c^2 + E_{\mathrm{kin}}$), and that the entropy increase matches the
total information destroyed.

\[
  A_{(ij),k} \;=\;
  \begin{cases}
    m_i + m_j               & k=1 \quad\text{(total annihilated mass)} \\
    2(m_i + m_j)c^2 + E_{\mathrm{kin}} & k=2 \quad\text{(energy released, eV)} \\
    t_{\mathrm{annihil}}    & k=3 \quad\text{(simulation time)} \\
    \rvec_{\mathrm{annihil}} & k=4 \quad\text{(position)} \\
    \Rec_i + \Rec_j          & k=5 \quad\text{(total information destroyed)} \\
    \Delta S_{\mathrm{annihil}} & k=6 \quad\text{(entropy produced)} \\
  \end{cases}
\]

\subsection{Annihilation Guard Conditions}

Annihilation is triggered when:
\begin{enumerate}[leftmargin=1.5em]
  \item The pair $(i, j)$ are confirmed particle-antiparticle.
  \item $d_{ij} \leq r_{\mathrm{annihil}}$ (default $10^{-4}\ang$ for
    leptons, $10^{-5}\ang$ for hadrons).
  \item $E_{\mathrm{kin}}^{\mathrm{rel}} \leq E_{\mathrm{annihil}}^{\mathrm{max}}$.
\end{enumerate}

\begin{smlaw}
\textbf{First Law of Identity} at annihilation: the information destroyed,
$\Rec_i + \Rec_j$, is entirely unrecoverable.  The engine writes the full
$\Rec$ budget to the entropy ledger and increments the global annihilation
entropy counter $S_{\mathrm{annihil}}^{\mathrm{glob}}$.  Annihilation without
entropy accounting raises \texttt{ERR\_ANNIHIL\_INFO\_MISSING}.
\end{smlaw}

%%═══════════════════════════════════════════════════════════════════════════════
\section{Detector Projection and Reconstruction}\label{sec:detector}
%%═══════════════════════════════════════════════════════════════════════════════

The detector projection module is the formal bridge between simulation and
experiment.  It mirrors the actual reconstruction pipeline used at particle
colliders: raw detector hits are processed through clustering, track-seeding,
momentum fitting, and vertex-finding to produce a list of reconstructed
particles with four-momenta and identity assignments.

In a real experiment this pipeline runs in the opposite direction to the
simulation: measured hits are the inputs and reconstructed particles are the
outputs.  Here the pipeline runs forward: true particle states are the inputs
and synthetic hits are generated, then processed by the same reconstruction
algorithm.  This allows direct comparison between truth-level identity
(from $\Omega(t)$) and reconstructed identity ($\tilde{\Pid}$), which is the
key figure of merit for the identity formalism: if the reconstruction fails
to separate two particles that the formalism treats as distinguishable, their
$\Rec$ values are too close and should be revised.

The detector module therefore acts as a continuous test of the information
accounting: it tells you whether the $\Rec$ values you have assigned to your
particles are experimentally meaningful or not.

\subsection{Projection Operator}

A detector $D_\alpha$ with acceptance function $A_\alpha(\Pid_i)$ produces
the hit vector:
\[
  \mathbf{h}^{(\alpha)}(t) \;=\;
  \sum_i A_\alpha(\Pid_i) \,
  \boldsymbol{\delta}\!\left(\rvec_i - \rvec_{D_\alpha}\right)
  \otimes \mathcal{E}_i(t) ,
\]
where $\otimes$ denotes the detector response convolution.

\subsection{Reconstruction Algorithm}

\begin{enumerate}[leftmargin=1.5em]
  \item \textbf{Hit clustering:} Group hits within spatial resolution
    $\delta x_{\mathrm{det}}$ and time window $\delta t_{\mathrm{det}}$.
  \item \textbf{Track seeding:} Each cluster with $n_{\mathrm{hit}} \geq 3$
    seeds a candidate track.
  \item \textbf{Kalman filter:} Propagate track hypothesis through the
    material budget and fit momentum.
  \item \textbf{Vertex finding:} Intersect compatible tracks at the primary
    vertex; record secondary vertices for decay topologies.
  \item \textbf{Identity assignment:} Assign reconstructed identity
    $\tilde{\Pid}$ and compare to truth-level $\Pid$ from simulation state.
\end{enumerate}

\begin{lstlisting}
# detector reconstruction output (.reco)
STEP   TRACK_ID  PX_GEV  PY_GEV  PZ_GEV  MASS_GEV  PID_TRUE  PID_RECO
50000  0001      0.4523  0.1204  2.3311   0.1396    211       211
50000  0002      -0.3901 0.0891  1.9823   0.1396    -211      -211
\end{lstlisting}

%%═══════════════════════════════════════════════════════════════════════════════
\section{Solver Comparison State}\label{sec:solver}
%%═══════════════════════════════════════════════════════════════════════════════

The solver comparison state introduces a new dimension to numerical
validation.  Classical MD validation typically checks energy conservation
and trajectory divergence between different integrators.  UIISCF adds
a third metric: information divergence $\delta\Rec_{\alpha\beta}$.

The significance of this metric is that it reveals a class of solver pathology
that energy conservation alone cannot detect.  A non-symplectic integrator
(Runge-Kutta 4, Adams-Bashforth) can conserve energy to high precision on
short runs while silently misattributing information flow between particle
pairs.  The $\delta\Rec$ metric catches this because the information-flow
equation is sensitive to the order of force evaluation in a way that the
energy integral is not.  Symplectic integrators (velocity Verlet, PEFRL,
Yoshida 6) preserve the information budget automatically because they preserve
the symplectic structure of the phase space; non-symplectic integrators do not.

Two or more solver threads share the same initial state $\Omega(0)$ and evolve
independently.  Comparison at checkpoints:
\[
  \delta\Rec_{\alpha\beta}(t_c) \;=\;
  \left|
    \sum_i \Rec_i^{(\alpha)}(t_c) - \sum_i \Rec_i^{(\beta)}(t_c)
  \right| .
\]

\begin{table}[H]
\centering
\caption{Available integrators for solver comparison.}
\label{tab:solvers}
\begin{tabular}{@{}llll@{}}
\toprule
Integrator & Code & Order & Symplectic \\
\midrule
Velocity Verlet  & \texttt{vv}    & 2 & Yes \\
Leapfrog         & \texttt{lf}    & 2 & Yes \\
Runge-Kutta 4    & \texttt{rk4}   & 4 & No \\
PEFRL            & \texttt{pefrl} & 4 & Yes \\
Yoshida 6        & \texttt{y6}    & 6 & Yes \\
Adams-Bashforth  & \texttt{ab4}   & 4 & No \\
\bottomrule
\end{tabular}
\end{table}

\begin{smlaw}
\textbf{First Law of Identity} provides a solver-independent invariant:
$\delta\Rec_{\alpha\beta}$ grows monotonically for non-symplectic solvers
due to artificial information injection, while symplectic solvers maintain
bounded $\delta\Rec$.  A symplectic solver showing unbounded $\delta\Rec$
growth indicates a force-field bug, not a solver defect.
\end{smlaw}

\begin{lstlisting}[language={}, caption={Solver comparison configuration block.}]
[solver_compare]
enabled             = true
solvers             = ["vv","pefrl"]
checkpoint_interval = 1000
compare_metrics     = ["pos_rmsd","energy_drift","info_divergence"]
output_path         = "solver_cmp.csv"
\end{lstlisting}

%%═══════════════════════════════════════════════════════════════════════════════
\section{Slimmed-Down Runtime Block}\label{sec:runtime}
%%═══════════════════════════════════════════════════════════════════════════════

Full UIISCF bookkeeping has a cost.  The per-particle $\Rec$ ODE adds one
float per particle per step; the sparse conductance matrix adds memory
proportional to the number of proximate pairs; the event log grows with
simulation time.  For large-scale production runs --- millions of atoms,
millions of steps --- this cost may be unacceptable.

The slim mode addresses this by suspending all information-theoretic
bookkeeping while preserving the classical conservation checks.  The
fundamental invariants (energy, momentum, charge, particle count) are still
validated.  The identity-information block is frozen: $\Rec_i$ and
$\Cond_{ij}$ are not updated.  Event logging is reduced to a summary.
The result is a run that is indistinguishable from a classical MD run in
terms of output, but which can be trivially upgraded to full UIISCF mode
by changing one configuration line.

\begin{table}[H]
\centering
\caption{Feature availability in full vs.\ slim runtime modes.}
\label{tab:slim}
\begin{tabular}{@{}lcc@{}}
\toprule
Feature & Full UIISCF & Slim mode \\
\midrule
Positions, velocities, forces  & \checkmark & \checkmark \\
Kinetic / potential energy      & \checkmark & \checkmark \\
Thermostat / barostat           & \checkmark & \checkmark \\
Per-atom $\Rec_i$ tracking      & \checkmark & --- \\
Conductance matrix $\Cond_{ij}$ & \checkmark & --- \\
Scale quota $\Quota_S$          & \checkmark & Aggregate only \\
Event log (\texttt{.evlog})     & \checkmark & Summary only \\
Subatomic identity matrices     & \checkmark & --- \\
Solver comparison               & \checkmark & --- \\
Dark-sector projection          & \checkmark & --- \\
Detector reconstruction         & \checkmark & --- \\
\bottomrule
\end{tabular}
\end{table}

%%═══════════════════════════════════════════════════════════════════════════════
\section{Version Constraints}\label{sec:version}
%%═══════════════════════════════════════════════════════════════════════════════

The UIISCF is implemented in stages, gated by version.  This is a deliberate
engineering decision: delivering the full formalism at once would require
users to adopt the entire information-tracking infrastructure before they
have any practical need for it.  Version gating allows the formalism to be
introduced incrementally, with each version boundary corresponding to a
tested, validated capability plateau.

Version 5.1.4 is the classical MD baseline.  Everything up to and including
the matrix-force dispatch table, event logging, and formation-defect tracking
is available without any subatomic content.  This version is suitable for
the vast majority of chemistry and materials-science simulations.

Version 5.1.5 activates the subatomic layer: quarks, leptons, bosons,
information tracking, annihilation, and the full UIISCF state vector.  It
requires explicit opt-in for every subatomic feature and will reject
configurations that enable subatomic channels without the information tracking
infrastructure.

\begin{table}[H]
\centering
\caption{Feature version gate table.}
\label{tab:version}
\begin{tabular}{@{}llp{7.2cm}@{}}
\toprule
Feature & Min.\ version & Notes \\
\midrule
Classical MD, matrix-force policy  & v5.1.4 & Baseline \\
Event logging (\texttt{.evlog})    & v5.1.4 & Bond, formation, defect events \\
Formation-defect bridge            & v5.1.4 & Formation trajectory format required \\
Per-atom $\Rec_i$ tracking         & v5.1.5 & \texttt{info\_track = true} \\
Conductance matrix $\Cond_{ij}$    & v5.1.5 & Sparse; max $10^6$ non-zeros \\
Scale identity ladder (subatomic)  & v5.1.5 & \texttt{subatomic\_mode = true} \\
Quark/lepton identity matrices     & v5.1.5 & Requires \texttt{full\_qcd = false} \\
Boson identity matrix              & v5.1.5 & Photon proxy + weak proxy \\
Annihilation events                & v5.1.5 & Lepton and hadron channels \\
Solver comparison state            & v5.1.5 & Any two registered integrators \\
Dark-sector projection             & v5.1.5 & Phenomenological only \\
Detector projection/reconstruction & v5.1.5 & Synthetic; no hardware interface \\
Proper-time correction             & v5.2   & \textit{Planned} \\
Explicit electron tracking         & v5.2   & \textit{Planned} \\
Full relativistic boost            & v5.2   & \textit{Planned} \\
\bottomrule
\end{tabular}
\end{table}

\subsection{Activation Checklist (v5.1.5)}

Before enabling any subatomic feature, verify:
\begin{enumerate}[leftmargin=1.5em]
  \item \texttt{subatomic\_mode = true} in the kernel block.
  \item \texttt{info\_track = true} in the run or kernel block.
  \item \texttt{full\_qcd = false} in any quark or boson block (no default).
  \item Kernel version declared at \texttt{"5.1.5"} or higher.
  \item Event log path configured.
\end{enumerate}

\begin{smlaw}
\textbf{Second Law of Identity} applies at the version level: a configuration
with $m \geq 3$ concurrent v5.1.5 feature blocks enabled and all guard
conditions satisfied will automatically activate the full UIISCF state.
A partial configuration raises \texttt{ERR\_PARTIAL\_UIISCF} and is rejected.
Partial activation is not supported; the formalism is all-or-nothing at v5.1.5.
\end{smlaw}

%%═══════════════════════════════════════════════════════════════════════════════
\section{Conservation Summary and Global Invariants}\label{sec:conservation}
%%═══════════════════════════════════════════════════════════════════════════════

\begin{table}[H]
\centering
\caption{Global invariants in UIISCF.}
\label{tab:invariants}
\begin{tabular}{@{}lllp{5.5cm}@{}}
\toprule
Quantity & Symbol & Conservation type & Violation handler \\
\midrule
Total energy       & $E_{\mathrm{tot}}$ & Approximate (NVE) & Rescale thermostat \\
Total momentum     & $\mathbf{P}$       & Exact             & Fatal error \\
Total charge       & $Q_{\mathrm{tot}}$ & Exact             & Fatal error \\
Particle count     & $N$                & Exact (non-annihil)& Fatal error \\
Lepton number      & $L_{\mathrm{tot}}$ & Exact             & Fatal error \\
Baryon number      & $B_{\mathrm{tot}}$ & Exact             & Fatal error \\
Total $\Rec$       & $\sum_i \Rec_i$    & Monotone non-increasing & Logged; not fatal \\
Entropy            & $S_{\mathrm{tot}}$ & Monotone non-decreasing & Logged; not fatal \\
Scale quota sum    & $\sum_S \Quota_S$  & Bounded change per step & Warning if $|\Delta\Quota_S| > 10$ \\
\bottomrule
\end{tabular}
\end{table}

\noindent
The information invariants (rows 7--8) are soft: they are tracked and logged
but do not cause fatal errors.  This is intentional.  The First Law of
Identity is a physical law, not an implementation constraint, and apparent
violations indicate interesting physics rather than bugs.  Only a spontaneous
\emph{increase} in $\sum_i \Rec_i$ without a corresponding event is treated
as a bug (\texttt{ERR\_INFO\_CREATED}).

%%═══════════════════════════════════════════════════════════════════════════════
\section{Complete Reference Script}\label{sec:refblock}
%%═══════════════════════════════════════════════════════════════════════════════

A complete annotated configuration script demonstrating all UIISCF features:

\begin{lstlisting}[language={}, caption={Complete UIISCF reference configuration (v5.1.5).}]
[meta]
version        = "5.1.5"
script_id      = "UIISCF-REF-001"
description    = "Full UIISCF feature demonstration"

[material]
file           = "system.xyz"
format         = "xyz"

[run]
integrator     = "velocity_verlet"
dt             = 0.5
steps          = 500000
thermostat     = "nose_hoover"
thermostat_T   = 300.0
barostat       = "parrinello_rahman"
barostat_P     = 1.0
info_track     = true
info_dissip    = 0.001

[kernel]
mode                = "full"
subatomic_mode      = true
event_rate_density  = 0.01
event_energy_guard  = 5.0
event_queue_depth   = 1000
event_log_path      = "run.evlog"
energy_tol          = 1.0e-4
momentum_tol        = 1.0e-6

[quark]
enabled         = true
flavours        = ["u","d","s"]
alpha_s         = 0.18
kappa_string    = 0.18
full_qcd        = false          # MANDATORY
cai_channel     = true

[boson]
enabled         = true
photon_proxy    = true
weak_proxy      = false
full_qcd        = false          # MANDATORY

[dark_sector]
enabled         = false

[solver_compare]
enabled             = true
solvers             = ["vv","pefrl"]
checkpoint_interval = 5000
output_path         = "solver_cmp.csv"

[observe]
interval        = 100
properties      = ["temperature","pressure","energy","Xi_total","quota_S"]
output          = "observables.chart"

[export]
trajectory          = "run.xyza"
formation           = "run.xyzf"
checkpoint          = "run.xyzc"
checkpoint_interval = 10000
reco_output         = "run.reco"
\end{lstlisting}

%%═══════════════════════════════════════════════════════════════════════════════
\section*{Summary of the Two Laws}
\addcontentsline{toc}{section}{Summary of the Two Laws}
%%═══════════════════════════════════════════════════════════════════════════════

\begin{smlaw}
\textbf{First Law of Identity --- Conservation of Distinguishability (canonical form).}
\[
  \Pid_i \leftrightarrow \Pid_j
  \;\implies\;
  \Delta\Rec_{ij} < 0 \;\;\text{or}\;\; \Delta S_{ij} > 0 .
\]
No physical interaction preserves or increases total recoverable
distinguishability.  Information is either converted to entropy or lost to
the environment.  This law holds at all scales and is scale-independent.
It is the reason detector reconstruction is a lossy process; it is the reason
annihilation destroys identity irreversibly; it is the reason that every
interaction in any simulation must reduce the information budget of at least
one participant.
\end{smlaw}

\begin{smlaw}
\textbf{Second Law of Identity --- Scale-Coalescence Birth Criterion (canonical form).}
\[
  m \geq 3,\quad
  d_{ij}^{(n)} \leq \lbond^{(n)},\quad
  B_n\!\left(C_n\right) \geq \tau_B
  \;\implies\;
  \mathrm{Exist}\!\left(\Pid_*^{(n+1)}\right) = 1,\quad
  \Delta\Quota_S = +1 .
\]
Complexity emerges through coalescence.  Every stable higher-scale identity
requires at minimum three lower-scale identities satisfying proximity and
bond maturity.  The new identity inherits a reduced but positive information
budget, consistent with the First Law.  Quarks become hadrons; nucleons become
nuclei; atoms become molecules.  The universe does this at every scale,
continuously, and has done so for $1.4 \times 10^{10}$ years before anyone
wrote it down.
\end{smlaw}

%%─────────────────────────────────────────────────────────────────────────────
%% Colophon
%%─────────────────────────────────────────────────────────────────────────────
\vspace{2em}
\noindent\rule{\textwidth}{0.6pt}

\begin{center}
\small\textit{%
  What CERN produces in abundance is not particles but information about
  interactions.\\[3pt]
  Working backwards from that information to a coherent picture of identity
  is the whole problem.\\[4pt]
  Theoretical Division \quad\textbar\quad
  UIISCF Rev.~1.0 \quad\textbar\quad \today
}
\end{center}

\end{document}
